% Chapter 2 — The fit map as a random dynamical system
% G2 content: PO-8 solved (closed forms + numeric certification); PO-9..12 framed with pointers.
\documentclass[11pt]{article}
\usepackage{amsmath,amssymb,amsthm}
\newtheorem{theorem}{Theorem}[section]
\newtheorem{proposition}[theorem]{Proposition}
\newtheorem{lemma}[theorem]{Lemma}
\newtheorem{corollary}[theorem]{Corollary}
\newtheorem{conjecture}[theorem]{Conjecture}
\newtheorem{definition}[theorem]{Definition}
\newtheorem{remark}[theorem]{Remark}
\newcommand{\R}{\mathbb{R}}
\title{Chapter 2: The fit map as a random dynamical system}
\begin{document}
\maketitle

\section{The fit map}
Let $F:(y,\theta_0,\xi)\mapsto\theta_T$ denote the fitting procedure mapping a target signal $y$,
an initialization $\theta_0$, and optimization noise $\xi$ to final weights. The program's
protocols are interventions on $F$'s arguments: P-shared-det fixes $(\theta_0,\xi)$;
P-shared-stoch fixes $\theta_0$; P-random draws both; P-random-K draws $K$ independent
$(\theta_0,\xi)$ per $y$. Chapter goal: a quantitative anatomy of the nuisance structure
$F$ injects, beginning with a completely solvable microcosm.

\section{The one-neuron microcosm (PO-8) [PROVED forms; numeric certification]}

Model $g_{w,b,u,c}(t)=u\sin(wt+b)+c$ on $t\in[-1,1]$; target $y(t)=A\sin(\omega t+\varphi)+c_0$.
Population loss $\mathcal L(w,b,u,c)=\tfrac12\int_{-1}^{1}(g-y)^2\,dt$.

\begin{lemma}[Profiling out the linear parameters]\label{lem:profile}
For fixed $(w,b)$ with $s(t)=\sin(wt+b)$, the optimal $(u,c)$ solve the $2\times2$ normal
equations with Gram entries
\[
\langle s,s\rangle = 1-\frac{\cos 2b\,\sin 2w}{2w},\qquad
\langle s,1\rangle = \frac{2\sin b\sin w}{w},\qquad
\langle 1,1\rangle = 2,
\]
and target moments ($\mathrm{sinc}\,x:=\sin x/x$, value $1$ at $0$)
\[
\langle s,y\rangle = A\big[\cos(b-\varphi)\,\mathrm{sinc}(w-\omega)-\cos(b+\varphi)\,
\mathrm{sinc}(w+\omega)\big] + c_0\langle s,1\rangle,\qquad
\langle 1,y\rangle = 2A\sin\varphi\,\mathrm{sinc}\,\omega + 2c_0 .
\]
The profiled loss is
$\mathcal L^*(w,b)=\tfrac12\big(\langle y,y\rangle - r^\top \mathrm G^{-1} r\big)$ with
$\mathrm G=\begin{psmallmatrix}\langle s,s\rangle&\langle s,1\rangle\\ \langle s,1\rangle&2
\end{psmallmatrix}$, $r=\begin{psmallmatrix}\langle s,y\rangle\\ \langle 1,y\rangle
\end{psmallmatrix}$, defined wherever $\det \mathrm G\neq0$ (all $(w,b)$ with $s\not\equiv$
constant).
\end{lemma}

\begin{proof} Elementary integrals via product-to-sum identities; least squares in $(u,c)$ is
linear. \end{proof}

\begin{proposition}[Global minima are exactly the symmetry copies]\label{prop:zeroset}
Assume $A\ne0$, $\omega\ne0$. Then $\mathcal L^*(w,b)=0$ iff
$(w,b)\in\{(\varepsilon\omega,\ \varepsilon\varphi+\pi k):\varepsilon\in\{\pm1\},k\in\Z\}$ —
precisely the $(w,b)$-projection of the $D_\infty$-orbit of the realizing parameters (with the
matching $u=(-1)^k\varepsilon A$, $c=c_0$).
\end{proposition}

\begin{proof}
$\mathcal L^*=0$ iff $y\in\mathrm{span}\{s,1\}$, i.e. $A\sin(\omega t+\varphi)=u\sin(wt+b)+
(c-c_0)$ on $[-1,1]$. Both sides extend to entire functions, so equality holds on $\R$; by
Lemma 1.6.2 of Ch.~1 (uniqueness for finite sine sums), the frequency sets $\{|\omega|\}$ and
$\{|w|\}$ coincide with matched complex amplitudes, forcing $w=\varepsilon\omega$,
$ue^{\mathrm i b}$ matched to $\varepsilon$-normalized $Ae^{\mathrm i\varphi}$ up to the
$\rho/\tau$ ambiguity, and $c=c_0$. These are exactly the displayed copies.
\end{proof}

\begin{remark}[Spurious minima and basin census — numeric certification]
$\mathcal L^*$ additionally possesses strictly positive local minima produced by the
finite-domain leakage kernels $\mathrm{sinc}(w\mp\omega)$ (sidelobes at spacing $\approx\pi$ in
$w$). These are certified numerically, not asymptotically: \texttt{scripts/02\_microcosm\_po8.py}
(i) evaluates the closed form on a fine grid and checks Proposition~\ref{prop:zeroset}'s zero set
to tolerance; (ii) locates all local minima on the grid with Hessian-PD confirmation;
(iii) runs descent from init ensembles and classifies endpoints into three types:
\emph{global} (an orbit copy), \emph{spurious} (a positive-loss sidelobe basin), and the
\emph{degenerate ridge} $w\approx0$, where the Gram matrix of $\{s,1\}$ degenerates and the
profiled loss equals the best-constant fit (pseudo-inverse semantics — the naive $G^{-1}$ formula
blows up there, a bug caught by the certification run itself); (iv) reports the census vs the
init range. Measured (target $\omega=7$, $n{=}900$ inits per range): range $\pm2$ — $100\%$ ridge
capture, the orbit is \emph{never} reached; $\pm5$ — $26\%$ global / $74\%$ ridge; $\pm10$ —
$62\%$ global (sweet spot near $|\omega|$); $\pm20$ — $51\%$ of inits captured by four spurious
sidelobe basins ($32\%$ global). Twenty isolated minima in the fundamental domain (one global,
nineteen spurious), all Hessian-checked on the grid. The census is thus \emph{non-monotone} in
the init range with three failure regimes, quantifying in closed form the folklore that sinusoidal
fitting requires frequency-matched initialization. Artifacts: \texttt{results/microcosm/} (F4).
This is the first place ``symmetry copies vs genuine basins (vs degenerate ridge)'' becomes a
computed, not rhetorical, distinction — the same discriminator logic S4 applies at scale.
\end{remark}

\section{Program for PO-9 to PO-12 [wired, empirical]}
\textbf{PO-9 (laziness).} Relative NTK movement and weight travel vs width; consequence: under
P-shared-det, $\varepsilon$-lazy fitting makes $F(\cdot,\theta_0)$ approximately affine in $y$,
predicting W1 linear-probe $\approx$ pixel linear-probe (registered direction P-D). External
support: shared-anchor linearity observed in arXiv:2607.04123; NTK conditioning of periodic INRs
in arXiv:2606.29414.
\textbf{PO-10 (basins).} $d_G(\theta,\theta')=\min_g\|\theta-g\theta'\|$ approximated by
\texttt{c\_align}; basin equivalence = small $d_G$ or post-alignment linear mode connectivity
(barrier $<\varepsilon_{\mathrm{barrier}}$, pre-registered in S4); estimator validated on
planted-$g$ pairs. Definitions to be reconciled with arXiv:2606.04754's effective function
classes at S4 prereg.
\textbf{PO-11 (chaos).} Finite-difference sensitivity of $F$ under P-shared-det; divergence-rate
curves; C-11 (non-monotone decodability vs fit length) registered at S4 prereg — delta over Papa
et al.'s RQ2 is the mechanism (rates, post-alignment dispersion, basin counts), not the phenomenon.
\textbf{PO-12 (sample complexity).} Finite-group benefits (Elesedy--Zaidi; Mei et al.; Bietti et
al.) adapted to $D_\infty$ via the truncated group $\{|j|\le J\}$ with $J$ from the measured 95\%
bias range $R_b$: effective orbit count $|G_{\mathrm{eff}}|\approx n!\,2^n(R_b/\pi)^n$ per layer,
with a truncation-error term to be stated when the S5 prereg fixes the estimator; predicted
ordering equivariant $\ge$ canonicalized $\ge$ augmented $\ge$ raw, gaps growing with
$|G_{\mathrm{eff}}|$ and shrinking with training-set size. The published failure of unbounded
phase augmentation (Shamsian et al., Table 3) is consistent with the truncation term blowing up
when augmentation samples $j$ far outside the data's bias range.

\end{document}
